\begin{table}[htbp]
  \caption{The fluency ceiling does not explain the sign of the primary
  test. If output overlap made a direction lexical enough to break fluency early,
  it would shorten the usable coefficient range and shrink an integrated area
  mechanically. It does not: overlap is unrelated to where the ceiling falls, and
  conditioning the primary test on the ceiling as well as on readability leaves
  the estimate unchanged.}
  \label{tab:ceiling}
  \centering
  \footnotesize
  \begin{tabular}{lcc}
    \toprule
    Quantity & Statistic & $p$ \\
    \midrule
    \multicolumn{3}{l}{\emph{Does output overlap truncate the sweep?}} \\
    Overlap vs.\ max usable $|\alpha|$ & Spearman $-0.096$ & $0.61$ \\
    Overlap vs.\ number of usable grid points & Spearman $0.006$ & $0.98$ \\
    Max usable $|\alpha|$ vs.\ controllability & Spearman $0.019$ & $0.92$ \\
    \addlinespace
    \multicolumn{3}{l}{\emph{Primary test under wider control sets}} \\
    Partial $\rho$ $|$ readability & $-0.450$ & \\
    Partial $\rho$ $|$ readability, max usable $|\alpha|$ & $-0.450$ & \\
    Partial $\rho$ $|$ readability, ceiling, $n$ usable & $-0.435$ & \\
    \bottomrule
  \end{tabular}
\end{table}
